\subsection{Task-conditioned computational erasure is not epistemic absence}
\label{sec:task-conditioned-erasure-authority}

A derived computational view may deliberately omit distinctions that have been certified as unnecessary for a registered quantity of interest without altering the canonical scientific state. Let $P$ be a content-bound source object and let
\[
Z=Q_{q,c}(P)
\]
be a task-conditioned derived view for quantity of interest $q$ and context $c$. The map $Q_{q,c}$ can change computational access---for example, which coordinates participate in retrieval, obstruction construction or downstream search---while the source object and its epistemic certificates remain unchanged.

The authority rule is therefore
\[
\text{coordinate absent from }Z
\not\Rightarrow
\text{coordinate false, disproved or scientifically irrelevant in }P.
\]
An erasure is local to the declared computational purpose. A coordinate suppressed for a stability question may remain essential for an intervention, fairness, mechanism or boundary question. The erasure ledger must therefore remain linked to the immutable source rather than replacing it.

\begin{proposition}[Derived-view authority non-escalation]
Suppose a task-conditioned view $Z=Q_{q,c}(P)$ is materialized only as a derived representation whose authority certificates are a subset of certificates reachable from the pinned source $P$. Then constructing, retrieving from or solving over $Z$ cannot by itself increase scientific authority of any claim in the canonical epistemic projection.
\end{proposition}

\begin{proof}
The quotient transition has codomain in derived computational state, not canonical scientific state. By construction it can inherit but cannot mint authority certificates, and the canonical source remains content-addressed and immutable. Hence any proposal or result obtained from $Z$ remains proposal-side until a separately licensed scientific update is supported by the ordinary evidence or verification path. In particular, computational suppression of a coordinate is not a scientific retraction, and success on the reduced view is not itself an authority certificate for the original problem.
\end{proof}

When a quotient-side computation produces a candidate solution, reconstruction and verification against the original source are distinct obligations. Thus
\[
\text{validated reduction}
\not\Rightarrow
\text{reduced solution correct}
\not\Rightarrow
\text{original scientific problem solved}.
\]
This is the same non-sovereignty principle used for workspaces, experience routing and semantic shortcuts: derived access structures can alter what the system computes next, but they do not silently rewrite what the scientific record is licensed to assert. Paper~II studies the structural mechanics and empirical usefulness of such task-conditioned reductions; the present paper uses them only to sharpen the authority boundary.
